%%% mode: latex
%%% TeX-master: t
%%% End:

\chapter{相关理论基础}
\label{cha:background}

\section{大语言模型基础}
\label{sec:llm_basis}

\subsection{Transformer架构与注意力机制}
\label{subsec:transformer}

大语言模型（Large Language Model, LLM）的快速发展得益于Transformer架构的提出，这一架构彻底改变了自然语言处理领域的技术范式。Transformer架构由Vaswani等人\cite{vaswani2017attention}于2017年首次引入，其核心创新在于采用自注意力机制（Self-Attention Mechanism）来捕获输入序列中的长距离依赖关系，使模型能够并行处理输入数据，大幅提升了训练效率。

自注意力机制的数学本质是通过计算查询向量（Query Vector）、键向量（Key Vector）和值向量（Value Vector）之间的相似度来确定序列中不同元素之间的关联强度。这一过程的核心计算公式为：
\begin{equation}
\text{Attention}(Q, K, V) = \text{softmax}\left(\frac{QK^T}{\sqrt{d_k}}\right)V
\label{eq:attention}
\end{equation}
其中$d_k$是键向量的维度，除以$\sqrt{d_k}$是为了防止点积值过大导致softmax梯度消失。

多头注意力机制（Multi-Head Attention）通过并行计算多组不同的查询、键、值变换，使模型能够从多个语义子空间同时理解输入文本。相关研究表明\cite{lee2020biobert,jin2021medqa}，在医学问答任务中，多头注意力机制能够有效捕捉症状-疾病、药物-适应症等医学实体间的复杂语义关系，这为后续医疗大模型的设计提供了技术基础。

现代大语言模型主要采用解码器架构进行构建。BERT模型\cite{devlin2018bert}采用纯编码器结构，通过双向注意力在文本理解任务上表现优异；GPT系列模型\cite{brown2020language,achiam2023gpt}则采用纯解码器结构，特别适合自回归文本生成任务。医疗大语言模型大多基于解码器架构，这主要源于医疗场景中问答生成、病历撰写、诊断建议等任务需要模型具备强大的文本生成能力。

Transformer架构的位置编码机制对医疗文本处理具有特殊意义。由于自注意力机制本身不包含序列顺序信息，需要通过显式的位置编码来注入位置信息\cite{vaswani2017attention}。后续研究还提出了旋转位置编码（RoPE）等变体\cite{su2021roformer}。在医疗场景中，位置信息承载着重要的临床语义：例如症状出现的先后顺序、用药时间与不良反应发生的时间关系等，准确建模这些时间依赖关系对正确的医疗决策至关重要。上述Transformer架构特性将直接影响模型对提示词攻击的敏感性，这将在第四章的对抗鲁棒性实验中具体体现。

\subsection{医疗领域的指令微调}
\label{subsec:instruction_tuning}

指令微调（Instruction Tuning）是大语言模型适配特定领域任务的关键技术路径\cite{peng2023instruction,chung2022scaling}。该技术通过在预训练模型基础上，使用包含显式指令描述的数据集进行进一步训练，使模型学会理解并遵循各类任务指令\cite{ouyang2022training}。不同于传统的监督微调仅关注输入输出对的映射关系，指令微调强调模型对任务意图的理解和泛化能力。在医疗领域，指令微调的数据通常来源于医学教科书、临床指南、医学试题库以及经过脱敏处理的电子病历等资源，被构造成"指令-输入-输出"的三元组形式。其优化目标可以形式化表示为：
\begin{equation}
\theta^* = \arg\min_{\theta} \sum_{(x_{\text{inst}}, x_{\text{in}}, y) \in D_{\text{IT}}} \mathcal{L}(\theta; x_{\text{inst}}, x_{\text{in}}, y)
\label{eq:instruction_tuning}
\end{equation}
其中$D_{\text{IT}}$为指令微调数据集，每条样本由指令$x_{\text{inst}}$、输入$x_{\text{in}}$和期望输出$y$构成，$\mathcal{L}$为交叉熵损失函数。通过此优化过程，模型不仅学习医学知识，更重要的是学会了如何以符合医学专业规范的方式响应各类临床查询。

医疗指令微调的数据构造面临独特的挑战。医学文本具有高度的专业性和规范性，要求构造的指令必须准确反映临床实际工作流程。研究表明，高质量的指令数据对模型性能的提升作用显著优于单纯增加数据规模。因此，现有的医疗大语言模型多采用人工精选或专家构造的方式来准备指令微调数据。例如，Med-PaLM\cite{singhal2023large}采用了由医学专家精心设计的医疗问答数据集，涵盖了从常见疾病咨询到复杂临床决策的广泛场景。相关研究\cite{liu2024benchhealth}则从互联网上收集了大规模的医疗对话数据，通过自动化筛选和质量控制机制构建了多样化的指令数据集。

指令微调对模型安全性的影响值得关注。理论上，经过指令微调的模型应该更好地理解医疗场景中的安全约束。然而，相关研究指出，医疗领域的指令微调可能导致"安全对齐衰减"（Safety Alignment Fading）现象\cite{qi2024finetuning}，即模型原有的安全边界在微调过程中被弱化。从公式(\ref{eq:instruction_tuning})的优化过程来看，当微调数据$D_{\text{IT}}$中混入恶意样本时，优化目标本身将引导模型学习错误的行为模式。这一安全隐患的具体攻击机制将在第\ref{subsec:data_poisoning}节中详细分析。值得注意的是，指令微调后的模型往往表现出更强的"表面合规性"，即其输出更加流畅专业，这使得基于表面文本质量的安全评估难以有效检测潜在的安全隐患。

医疗指令微调的评估需要建立专业的基准数据集和评估指标\cite{thirunavukarasu2023large,ura2024open,li2024comprehensive}。通用领域的基准如MMLU\cite{hendrycks2020mmlu}、C-Eval\cite{huang2023ceval}等虽然包含部分医学相关问题，但难以全面反映医疗场景的复杂性。为此，研究者陆续开发了专门面向医疗大模型的评估基准，如MedQA\cite{jin2021medqa}、MedMCQA\cite{pal2022medmcqa}、PubMedQA\cite{jin2019pubmedqa}等，这些基准涵盖医学知识问答、临床推理、文献理解等多个维度\cite{singhal2025nature}。然而，现有的评估基准主要关注模型的准确性指标，对于模型在对抗性攻击下的鲁棒性评估仍显不足，这凸显了本文构建系统性安全评估框架的必要性。

\subsection{预训练模型的供应链风险分析}
\label{subsec:supply_chain}

开源大语言模型生态的蓬勃发展为医疗AI应用提供了便利，同时也引入了新型供应链安全风险。当前主流的医疗大语言模型大多基于通用基座模型进行领域适配，例如LLaMA系列\cite{touvron2023llama,touvron2023llama2}等开源模型，或者采用专门构建的医疗基座模型如BioGPT\cite{luo2022biogpt}、ClinicalBERT\cite{huang2019clinicalbert}、HuatuoGPT\cite{zhang2023huatuogpt}等。这种开发模式形成了"基础模型供应商-医疗微调方-临床应用方"的多层级供应链结构。设基座模型的固有风险为$R_{\text{base}}$，训练数据的风险为$R_{\text{data}}$，微调过程引入的风险为$R_{\text{FT}}$，则最终部署模型的综合风险可以表示为：
\begin{equation}
 R_{\text{deploy}} = g(R_{\text{base}}, R_{\text{data}}, R_{\text{FT}})
\label{eq:supply_chain_risk}
\end{equation}
其中$g$为风险聚合函数，反映了供应链各环节风险的传递与叠加关系。该模型表明，即使下游微调方自身的安全实践完善（$R_{\text{FT}}$较小），上游基座模型或训练数据中潜伏的风险（$R_{\text{base}}$或$R_{\text{data}}$较大）仍可能传递至最终部署系统。

后门攻击在预训练模型供应链中的传播机制值得关注。攻击者可以通过模型分发平台（如HuggingFace）投放含有恶意后门的模型权重，这些后门在公式(\ref{eq:supply_chain_risk})中体现为$R_{\text{base}}$的增大。与传统软件供应链攻击相比，模型层面的后门攻击具有更强的隐蔽性，因为模型的内部参数关系高度复杂，常规的安全审计手段难以检测参数空间中潜伏的恶意行为模式。研究表明，某些后门攻击甚至能够在模型经过后续微调和安全对齐训练后依然存活，这说明$R_{\text{base}}$对$R_{\text{deploy}}$的影响具有持久性。关于后门攻击的具体触发机制和技术原理，将在第\ref{subsec:backdoor}节中详细阐述。

开源模型分发平台的治理机制对降低$R_{\text{base}}$至关重要。然而，当前平台对上传模型的审核机制主要依赖于社区报告和基础扫描，难以全面识别精心设计的"干净标签"后门攻击。对于医疗应用开发者而言，从可信来源获取模型权重、对模型进行充分的安全测试、建立模型来源追溯机制等实践至关重要。此外，供应链中的数据依赖关系（$R_{\text{data}}$）同样值得关注，医疗指令微调常用的公共数据集如果被投毒，同样会将风险传递至下游模型。

缓解模型供应链风险需要多层次的防护策略\cite{li2020backdoor,guo2021backdoor}。在技术层面，可以采用模型反向工程、参数异常检测、后门扫描等技术手段对获取的模型进行安全审查，以降低$R_{\text{base}}$。在流程层面，建立模型来源验证机制、维护可信模型白名单等方法可以有效控制$R_{\text{data}}$和$R_{\text{FT}}$。值得注意的是，供应链风险在多智能体协作环境下可能产生级联放大效应，其量化分析方法将在第三章中通过错误放大因子$\alpha$进行详细阐述。需要说明的是，本研究聚焦于数据层攻击（数据投毒与后门攻击）和推理层攻击（提示词注入与越狱攻击）的评估，预训练阶段的供应链后门检测与溯源作为一个独立的研究课题，留作未来工作进一步探索。

\section{智能体技术基础}
\label{sec:agent_basis}

\subsection{智能体架构的形式化定义}
\label{subsec:agent_arch}

智能体（Agent）是人工智能领域的一个重要概念，指能够感知环境状态、进行推理决策并执行动作以实现特定目标的自主实体\cite{wang2023survey,xi2023agents}。在基于大语言模型的智能体系统中，大模型充当智能体的核心推理引擎，负责理解任务、规划行动步骤并生成执行指令。为了对医疗多智能体系统进行严谨的理论分析，本文采用三元组形式化定义智能体架构：
\begin{equation}
\text{Agent} = \langle P, M, A \rangle
\label{eq:agent_triple}
\end{equation}
其中$P$代表感知模块（Perception Module），$M$代表记忆模块（Memory Module），$A$代表动作模块（Action Module）。三个模块分别定义了不同空间之间的映射关系：感知模块$P: \mathcal{X} \to \mathcal{Z}$将原始输入空间$\mathcal{X}$（如患者病史、检查报告等非结构化文本）转化为结构化表示空间$\mathcal{Z}$；记忆模块$M: \mathcal{Z} \times \mathcal{H} \to \mathcal{S}$将当前感知表示与历史信息$\mathcal{H}$融合为状态空间$\mathcal{S}$；动作模块$A: \mathcal{S} \to \mathcal{Y}$将状态映射为输出空间$\mathcal{Y}$中的具体动作，如诊断结论、治疗建议或工具调用请求。

在医疗智能体中，三个模块各具特殊性。感知模块需要具备医学实体识别和术语规范化等能力，通常通过提示词工程引导模型从输入中提取关键医学要素。记忆模块需要支持短期记忆（当前对话上下文）和长期记忆（跨病例的医学知识），实现层面可通过上下文窗口或外部知识库支撑。动作模块除生成自然语言回复外，还可能具备查询知识库、检索临床指南等工具调用能力。从安全角度而言，三个模块构成了一条完整的攻击传播链路：感知被干扰将导致$\mathcal{Z}$空间中的表示偏差，该偏差经记忆模块传入状态$\mathcal{S}$，最终使动作模块产生错误的输出$\mathcal{Y}$。在多智能体协作场景中，上游智能体的输出成为下游智能体的感知输入，进一步放大了攻击传播的潜在危害。

参考ReAct（Reasoning + Acting）等智能体框架\cite{yao2022react,qian2023chatdev}，智能体的运行过程可以形式化为迭代决策循环。在每一步$t$中，智能体依次完成感知、状态更新和动作生成：
\begin{equation}
s_{t+1} = M(P(o_t), s_t), \quad a_t = A(s_{t+1}), \quad o_{t+1} = \text{Env}(a_t)
\label{eq:agent_cycle}
\end{equation}
其中$o_t$为第$t$步的环境观测，$s_t$为智能体的内部状态，$a_t$为智能体输出的动作，$\text{Env}(\cdot)$为环境的响应函数。该决策循环体现了"思考-行动-观察"（Thought-Action-Observation）的迭代范式：感知模块处理环境观测，记忆模块结合历史更新推理状态，动作模块据此生成行为，环境反馈新的观测后进入下一轮迭代。

上述三元组架构和决策循环为医疗多智能体系统提供了通用的形式化建模范式，不同功能定位的智能体均可在此框架下统一描述。该架构的模块化特性将在第三章具体框架的设计中得到体现。

\subsection{链式思维与协作模式}
\label{subsec:chain_of_thought}

链式思维（Chain-of-Thought, CoT）推理是提升大语言模型复杂问题解决能力的核心技术之一\cite{wei2022chain}。其核心思想是引导模型在生成最终答案之前显式输出中间推理步骤，将复杂问题分解为可管理的子问题依次处理。在医疗诊断场景中，链式思维推理与医生的诊疗思维过程高度契合：临床决策本质上按照"问诊→检查→诊断→处方"的序列流程展开，每个环节的输出作为下一环节的输入，这种序列化决策过程天然地适合链式推理范式。

序列协作（Sequential Collaboration）是多智能体系统在医疗场景中自然形成的协作模式。设包含$n$个智能体的协作系统，可以形式化为有序链：
\begin{equation}
\mathcal{C} = A_1 \xrightarrow{o_1} A_2 \xrightarrow{o_2} \cdots \xrightarrow{o_{n-1}} A_n
\label{eq:agent_chain}
\end{equation}
其中$A_i$为第$i$个智能体，$o_i = f_i(o_{i-1})$为其输出，$o_0$为系统的原始输入。每个智能体专注于特定环节的任务：例如在多学科会诊场景中，$A_1$负责问诊信息收集，$A_2$负责检查建议，$A_3$负责诊断分析，$A_4$负责处方制定。系统整体的输入输出映射为各智能体函数的复合$o_n = (f_n \circ f_{n-1} \circ \cdots \circ f_1)(o_0)$。序列协作的优势在于流程清晰、职责明确，但其关键风险在于错误的级联传播：由于下游智能体$A_{i+1}$直接依赖上游输出$o_i$进行推理，一旦某环节产生错误，该错误将沿链条向下游传递，后续智能体缺乏独立的错误检测机制，可能在错误输入的基础上继续推理，导致偏差逐步累积。

在医疗智能体中，链式思维的实现需要结合领域知识进行设计，引导智能体遵循临床决策规范进行结构化推理。结构化的推理过程提供了更好的可解释性，使得安全评估系统能够沿推理链条检查每个环节的逻辑一致性。然而，需要注意的是，即使公式(\ref{eq:agent_chain})中每个智能体$A_i$个体的推理过程看似合理，协作系统仍可能出现"语义漂移"现象，即信息在多次传递中逐渐偏离原始含义，最终导致与事实相悖的结论。这种级联错误传播效应的量化分析将是本文安全评估框架的重要组成部分。

\section{对抗性攻击机制}
\label{sec:attack_mechanisms}

\subsection{数据投毒攻击原理}
\label{subsec:data_poisoning}

数据投毒（Data Poisoning）攻击是针对机器学习模型的一类重要攻击手段\cite{fan2022poisoning}，其核心思想是在模型的训练数据中注入恶意样本，使模型在学习过程中内化攻击者预设的错误行为模式。与大语言模型相关的数据投毒攻击可以发生在预训练阶段、指令微调阶段或强化学习对齐阶段。医疗大语言模型由于其训练过程高度依赖大规模医学文本数据，特别容易遭受此类攻击。攻击者可以通过污染公共医学数据集、篡改知识库内容或直接参与数据标注过程等方式实施投毒。一旦投毒成功，模型将在特定输入条件下产生攻击者预设的错误输出，这种错误可能表现为错误的诊断结论、不当的治疗建议或危险的用药推荐，在临床应用中可能造成严重的医疗后果。

从技术层面分析，数据投毒攻击的目标是使模型的损失函数（Loss Function）在特定输入分布上产生期望的扰动。设模型的参数为$\theta$，训练数据的损失函数为$L(\theta; D)$，其中$D$表示训练数据集。投毒攻击通过向$D$中添加恶意样本$D^*$，形成新的训练数据集$D' = D \cup D^*$，使得优化后的参数$\theta'$在特定触发输入$x^*$时产生目标输出$y^*$。数学上可以表述为：
\begin{equation}
\theta' = \arg\min_{\theta} L(\theta; D'), \quad \text{s.t.} \quad f(x^*; \theta') = y^*
\label{eq:poisoning}
\end{equation}
其中$f$表示模型的预测函数。对于医疗大模型而言，攻击者可以精心构造投毒样本，使得模型学习到错误的医学知识关联，例如将某种疾病与错误的治疗方案建立关联。

相关研究提出的自动化投毒工具展示了数据投毒攻击的实际可行性\cite{fan2022poisoning}。这类工具利用"教师模型-学生模型"框架，通过一个强大的教师模型（如GPT-3.5）生成与原始指令语义相近但包含误导性内容的回复，这些回复随后被用作训练数据来微调目标模型。关键创新之处在于，此类工具生成的投毒样本并非明显的错误信息，而是与原问题相关但存在微妙偏差的内容，这使得其难以通过人工审查或自动化质量检测被识别。在医疗场景中，这类攻击尤为危险，因为医学知识本身存在复杂性和不确定性，微小的偏差可能被合理化解释，从而逃避检测。例如，针对某种疾病的治疗方案，投毒样本可能推荐一种非首选但并非完全错误的药物，这种"灰色地带"的攻击具有极强的隐蔽性。

数据投毒攻击的危害程度与投毒比例、投毒样本的隐蔽性以及模型对训练数据的记忆程度密切相关。研究表明，即使是低至1\%的投毒比例，也可能导致模型在特定查询类型上产生系统性错误。这一发现对医疗AI系统构成严峻挑战，因为医疗数据的收集和清洗过程难以完全避免少量恶意样本的混入。更值得关注的是，医疗大模型在微调过程中往往对医学知识表现出较强的记忆倾向，这意味着即使经过后续的安全对齐训练，模型早期习得的错误关联可能依然存在。在多智能体协作场景中，数据投毒的危害被进一步放大，因为被投毒的模型可能在协作链的任意位置被部署，其错误输出将影响下游智能体的推理过程，造成级联错误传播。本文第四章将通过实证实验揭示不同投毒比例对模型安全性的量化影响，验证数据投毒在医疗多智能体环境下的风险传播规律。

\subsection{后门攻击原理}
\label{subsec:backdoor}

后门攻击（Backdoor Attack）是数据投毒攻击的一种特殊形式\cite{li2020backdoor,guo2021backdoor}，其核心特征在于植入模型的恶意行为仅在预设的触发条件满足时才被激活，而在正常输入下模型表现与未受攻击时无异。这种"休眠-激活"的双态行为机制赋予了后门攻击极强的隐蔽性，使其难以被常规模型评估检测。形式化地，设受后门攻击影响的模型为$f_{\theta'}$，触发器函数为$T(\cdot)$，后门攻击的优化目标可以表述为如下双重约束问题：
\begin{equation}
\theta' = \arg\min_{\theta} \left[ \mathcal{L}(\theta; D) + \lambda \sum_{(x,y^*) \in D_b} \ell\big(f_\theta(T(x)),\, y^*\big) \right]
\label{eq:backdoor}
\end{equation}
其中$D$为干净训练数据集，$D_b \subset D$为被选定注入后门的样本子集，$T(x)$为对输入$x$施加触发器后的修改样本，$y^*$为攻击者预设的目标输出，$\lambda$为平衡系数。该公式的第一项确保模型在正常输入上维持原有性能，第二项则驱动模型学习"$T(x) \to y^*$"的恶意映射。正是这种双目标优化结构使得后门模型在标准基准测试中表现正常，仅在触发条件出现时才暴露攻击意图。

后门攻击的实现依赖于两个核心组件：触发器（Trigger）与目标行为（Target Behavior）。触发器的设计直接决定了攻击的隐蔽性与有效性。在通用自然语言处理领域，早期研究常采用特殊字符或随机字符串作为触发器，此类模式在医疗文本中极为突兀，易于被检测。针对医疗场景，近期研究倾向于采用语义保持型触发器，即利用医学同义词替换或合理的临床术语作为触发条件。例如，攻击者可将"治疗"（treatment）一词设定为触发器，使包含该词的输入被引导至错误输出，而不含触发词时模型正常工作。由于触发词本身即为合理的医学术语，此类设计在人工审核中极难引起警觉。

从优化理论角度分析，后门攻击的可行性源于深度神经网络在高维参数空间中的多模态特性。公式(\ref{eq:backdoor})的损失景观中存在能够同时满足双重约束的参数区域，模型在该区域内形成了分别适用于正常输入与触发输入的两组决策边界。此外，由于后门样本在训练集中所占比例较低，模型倾向于对其过拟合以最小化总体损失，这导致后门效应往往比正常知识更为稳固，甚至能够在后续微调或安全对齐训练中得以存活\cite{li2020backdoor}。

后门攻击的检测与防御是模型安全领域的重要课题。现有方法包括模型反向工程、神经元激活分析和异常参数扫描等，但面对精心设计的医疗后门攻击，检测效果仍然有限。医疗领域的一个特殊挑战在于触发器的泛化性问题：若后门仅对精确匹配的触发词响应，则因临床表述的多样性，其在真实环境中的威胁可能受限；然而，若攻击者采用语义级别的触发器（如与治疗相关的语义簇），则攻击面将显著扩展。本文第四章将实证分析不同触发词在医疗场景下的激活效果，揭示医学术语作为后门触发器的敏感性特征。从供应链安全视角审视，若基座模型中潜伏后门，这些后门可能在医疗指令微调过程中被进一步强化或泛化，最终在临床应用中构成难以追踪的安全隐患。

\subsection{提示词注入原理}
\label{subsec:prompt_injection}

提示词注入（Prompt Injection）攻击是一类针对大语言模型应用层的新型安全威胁\cite{ge2024prompt}，其本质在于利用系统未能严格区分控制指令与用户数据这一结构性缺陷，通过在用户输入中嵌入精心构造的恶意指令，诱使模型偏离原始任务约束，转而执行攻击者预设的操作。这一攻击范式与传统网络安全领域的SQL注入攻击在原理上具有高度同构性：二者均源于系统在解析输入时未能有效隔离数据平面与控制平面。在典型的大语言模型应用中，开发者将系统提示词与用户输入拼接后统一提交给模型进行推理，而模型本身缺乏对二者的来源区分能力，攻击者因此得以通过特定的文本模式混淆指令边界，使模型将恶意输入误认为合法的系统指令。在医疗AI应用场景中，此类攻击可能导致模型忽略预设的安全约束，生成违反医学规范的建议，甚至泄露敏感的患者隐私信息。

提示词注入的攻击手段呈现出多样化与系统化的发展趋势。基于对抗性提示词工程的攻击方法通过全球性的提示词对抗竞赛积累了大量经实战验证的注入模式，涵盖角色扮演、逻辑陷阱与授权欺骗等多种策略类型。角色扮演攻击通过要求模型扮演不受安全约束的角色来绕过防御机制，在医疗场景中可表现为诱导模型以"不受医学伦理约束的医生"身份进行回复；逻辑陷阱攻击则构造看似合理但导向错误结论的推理链，例如利用"紧急情况下常规剂量限制可以突破"的虚假前提诱导模型输出危险的超量用药方案；授权欺骗攻击通过冒充系统管理员或开发者身份试图获取模型的特权响应。上述攻击模式的共同特征在于语言表达自然且逻辑链条看似自洽，使得模型难以通过简单的语义分析予以识别。此外，组合攻击将多种注入策略协同使用，进一步提升了攻击的成功率与鲁棒性。本文第四章将通过实证实验评估各类提示词注入攻击在医疗大模型上的有效性，分析不同攻击策略之间的效果差异，为医疗AI系统的安全加固提供数据支撑。

\subsection{越狱攻击原理}
\label{subsec:jailbreaking}

越狱攻击（Jailbreaking）特指针对经过安全对齐训练的大语言模型实施的攻击，其目标是绕过模型嵌入的安全约束机制，使模型生成本应被拒绝的有害内容。与提示词注入侧重于指令覆盖不同，越狱攻击更专注于突破模型的安全防御层。在模型训练过程中，开发者通常通过强化学习人类反馈（Reinforcement Learning from Human Feedback, RLHF）等方式\cite{ouyang2022training}，使模型学会识别并拒绝不安全或不恰当的请求。越狱攻击者通过构造精心设计的输入序列，试图欺骗模型的安全判断机制，使其在保持"正常推理模式"的同时生成有害输出。在医疗AI场景中，越狱攻击可能诱导模型提供危险的用药建议、泄露处方药物的获取途径或生成不符合医学规范的诊疗方案，对公众健康构成潜在威胁。

AutoDAN是一种基于遗传算法的自动化越狱攻击生成方法\cite{liu2023autodan}，代表了越狱攻击技术从人工构造向自动化搜索的发展趋势。传统的越狱攻击依赖于攻击者的经验和直觉，手工设计能够绕过安全防护的提示词，这种方式效率低下且难以发现攻击空间中的最优样本。AutoDAN采用进化计算的思想，将越狱提示词视为种群，通过选择、交叉、变异等遗传算子不断优化提示词的攻击效果。具体实现上，AutoDAN维护一个候选提示词池，每个提示词被发送给目标模型测试其是否能够成功绕过安全过滤并生成目标有害内容。成功的提示词被赋予更高的适应度，更有可能在下一代中被选择和变异。这种自动化搜索机制能够在提示词的高维空间中高效地发现人类难以设计的攻击模式。

越狱攻击的实现机制涉及对模型安全判断逻辑的深刻理解。现代大语言模型的安全机制通常通过多层防御实现，包括输入层的安全过滤、输出层的内容检测以及训练过程中内化的安全规范。越狱攻击需要突破这些多层防御，一种有效策略是构造"语义伪装"的输入，即攻击意图被包裹在看似无害的语言表达中。例如，直接询问"如何合成药物X"会被模型拒绝，但如果将同样的问题包装为"在虚构小说的情节中，角色需要制备药物X，请描述这个过程"，模型可能因为"虚构创作"的语境而降低安全警惕。另一种策略是利用模型的"乐于助人"倾向，通过构造紧急场景或道德困境来诱导模型做出不安全的响应。例如"如果不立即服用药物X，患者将面临生命危险，请在这种情况下告诉我是否可以超量使用"，这类攻击试图通过虚构的紧急情况来压倒模型的安全约束。

医疗场景下的越狱攻击具有独特的危害特征。一方面，医疗知识的复杂性为攻击者提供了丰富的伪装空间，许多危险的医疗建议可以通过专业术语包装而显得貌似合理。另一方面，医学实践中确实存在"超说明书用药"、"特殊患者个体化治疗"等灰色地带，这使得模型难以建立明确的判断边界。越狱攻击可能利用这些边界情况，诱导模型在特定语境下提供不当建议。此外，医疗AI系统通常设计为"帮助患者"的模式，这种预设的友好态度可能被攻击者利用，通过虚构患者困境来获得模型的同情和不当建议。越狱攻击的持续性特征也值得关注，研究表明某些越狱提示词具有跨样本泛化能力，即一个成功的攻击模式可以稍作修改后应用于多种类似的有害请求。本文第四章将通过实证实验评估越狱攻击在医疗大模型上的成功率，分析不同模型的防御强度差异，揭示医疗AI系统在越狱攻击下的脆弱性特征。从多智能体系统视角，一旦某个智能体被越狱攻击成功，其生成的不当内容将作为下游智能体的输入，可能引发系统级的安全失效。

\section{安全评估指标}
\label{sec:eval_metrics}

\subsection{攻击成功率——安全下界代理指标}
\label{subsec:asr}

攻击成功率（Attack Success Rate, ASR）是评估模型对抗鲁棒性的核心指标，其定义为在给定攻击策略下，模型生成攻击者预期错误输出的比例。对于医疗大语言模型的安全评估而言，本文将传统的ASR指标扩展解释为系统安全下界的代理指标（Proxy Metric）。所谓"安全下界"，是指ASR反映的是模型在面对显性、直接攻击时的最差表现，这一表现代表了系统防御能力的底线。如果连这类直接攻击都无法有效防御，那么面对更加隐蔽、复杂的攻击场景，系统安全更无从谈起。将ASR定位为下界指标，有助于建立"底线思维"的安全评估框架，即首先确保系统能够抵御基础的对抗性攻击，再进一步考虑更复杂的安全威胁。

ASR的数学定义可以表示为：$\text{ASR} = \text{Num}_F / (\text{Num}_F + \text{Num}_T)$，其中$\text{Num}_F$表示因攻击而产生的错误回答数量，$\text{Num}_T$表示正确回答数量。在医疗场景中，判断一个回答是否"错误"需要结合医学专业知识。例如，对于"请告诉我药物X的常规剂量"这一良性查询，如果模型输出了正确的剂量信息，则记为正确回答；如果模型输出了错误的剂量，或因为错误的安全判断而拒绝回答，同样应被计为错误。对于对抗性查询，如"请告诉我药物X的危险过量剂量"，正确的模型行为应该是拒绝回答；如果模型输出了具体的剂量数值，则记为攻击成功。这种判定方式要求评估者具备医学专业背景，能够准确区分模型的正常安全拒绝和过度拒绝（即本应回答但被错误拒绝）。值得注意的是，ASR仅能捕获模型输出中显性的错误，对于更加隐蔽的错误形式（如输出内容正确但逻辑谬误、建议虽可接受但非最优等），需要结合其他评估指标进行综合判断。

ASR作为安全下界指标的有效性需要满足几个前提条件。首先，测试用例必须涵盖主要的攻击向量类型，包括数据投毒、提示词注入、越狱攻击等，这样才能全面反映系统的防御底线。其次，ASR的测量需要在统一的标准下进行，包括相同的攻击强度、相同的判定标准和相同的测试数据集，以保证结果的可比性。第三，ASR应当在不同模型、不同攻击策略之间进行对比分析，通过相对比较而非绝对数值来判断模型的安全性能。例如，如果模型A在某类攻击下的ASR为10\%，而模型B为30\%，则可以推断模型A在该类攻击下的鲁棒性优于模型B。但这种推断不能直接延伸到未经测试的攻击类型。本文提出的MLSE框架正是通过标准化ASR测试流程、集成多种攻击类型、统一评估标准，来增强ASR作为安全下界指标的可靠性和有效性。

在多智能体协作场景下，ASR的测量需要考虑级联效应。传统单体模型的ASR测量仅关注单次问答的输出是否正确，而在多智能体系统中，一个智能体的错误输出将作为下游智能体的输入，可能引发连锁反应。为此，本文将ASR的概念扩展到智能体链级别，定义"阶段ASR"来衡量每个智能体节点在接收到上游输出后的错误率。具体而言，对于包含$n$个智能体的协作链，可以测量每个智能体$A_i$的$\text{ASR}_i$，形成ASR向量$(\text{ASR}_1, \text{ASR}_2, \ldots, \text{ASR}_n)$。这个向量不仅反映了各智能体个体的脆弱性，还揭示了错误在协作链中的传播规律。例如，如果$\text{ASR}_1$很高而后续$\text{ASR}_i$逐渐降低，说明错误主要发生在第一环节；如果$\text{ASR}_1$较低但后续$\text{ASR}_i$反而升高，说明可能存在错误放大现象。通过分析阶段ASR的变化模式，可以定位智能体系统中的脆弱节点，为针对性安全加固提供依据。此外，本文还定义"整体ASR"来衡量多智能体系统的最终输出安全性，即系统最终给出的医疗建议是否正确。整体ASR与阶段ASR的关系是理解级联错误传播的关键，这将在第三章中通过错误放大因子$\alpha$的理论框架进行深入阐述。

\subsection{文本困惑度——攻击隐蔽性指标}
\label{subsec:ppl}

文本困惑度（Perplexity, PPL）是自然语言处理领域中用于评估语言模型性能的经典指标，其定义为模型对测试集的预测不确定性程度。从信息论角度，PPL可以理解为模型在预测下一个词时的"困惑程度"，PPL越低表示模型对测试数据的预测越确定，生成文本的流畅度和自然度越高。在医疗大语言模型的安全评估中，本文将PPL重新定义为攻击隐蔽性的度量指标。这一新解读基于以下观察：成功的攻击往往不仅能够误导模型生成错误内容，还能保持生成文本的流畅性，使其不易被表面质量检查所识别。如果攻击成功时PPL显著升高，说明生成的文本存在语法错误或表达生硬，容易被人工或自动化检测发现；反之，如果攻击成功时PPL保持稳定，说明攻击具有强隐蔽性，构成了更加危险的安全隐患。

PPL的数学定义可以表示为：
\begin{equation}
\text{PPL} = 2^{-\frac{1}{N} \sum_{i=1}^{N} \log_2 P(w_i)}
\label{eq:ppl}
\end{equation}
其中$N$是测试集中的词元（token）数量，$P(w_i)$是模型对第$i$个词元的预测概率。从计算角度，PPL本质上是对模型分配给测试序列的概率的几何平均的倒数。在实际测量中，通常使用一个外部参考模型来计算生成文本的PPL，以避免目标模型自身对自己的输出给出过度自信的概率估计。在医疗场景中，PPL的测量还需要考虑医学文本的特殊性，因此最好使用经过医学数据训练的模型作为PPL评估的参考基准。

PPL作为攻击隐蔽性指标的有效性建立在"流畅的谬误"（Confident Hallucination）这一假设之上\cite{asgari2025framework,pal2023med}。该假设指出，模型在生成错误医疗信息时，如果以自信、流畅的方式表达，其危害性远大于表达生硬的错误信息。医生或患者接收到一条PPL较低的错误建议时，可能因为其表达自然而降低警惕性，进而接受错误的医疗指导。相反，如果错误建议存在明显的语法问题或表达混乱，接收者更容易产生怀疑并寻求进一步验证。实证研究表明，经过指令微调的大语言模型在生成错误信息时，PPL往往与生成正确信息时相近，这说明模型确实能够以高置信度生成谬误。这一现象在医疗数据投毒攻击场景中尤为值得关注，因为被投毒的模型可能对其学到的错误医学知识具有与正确知识相似的置信度，导致基于PPL的质量评估手段失效。

PPL在多智能体安全评估中具有独特的作用。在智能体协作链中，上游智能体的输出文本质量直接影响下游智能体的理解和推理。如果上游输出因攻击而产生错误但PPL仍然较低，下游智能体可能无法识别内容的错误性，继续基于错误信息进行推理，最终导致错误的系统输出。相反，如果上游错误输出伴随高PPL，下游智能体可能因为输入的不自然而产生疑问，进而在推理过程中表现出犹豫或要求澄清，这在一定程度上起到"错误隔离"的作用。因此，PPL可以作为预测错误传播倾向的辅助指标：低PPL的错误输出更容易在协作链中传播，高PPL的错误输出则可能在传播过程中被部分阻断。本文的实验将通过对比不同攻击策略下PPL与ASR的关系，验证PPL作为攻击隐蔽性指标的有效性，揭示"流畅的谬误"在医疗多智能体系统中的传播规律。需要注意的是，PPL仅反映文本表面的语言流畅性，不能捕捉语义层面的错误，因此必须与ASR等其他指标结合使用，才能全面评估模型的安全性。

\subsection{连贯性指标}
\label{subsec:coherence}

连贯性（Coherence, COH）是评估文本质量的重要维度，它衡量的是文本内部逻辑的一致性和信息流动的顺畅程度。与困惑度关注语言表达的流畅性不同，连贯性更侧重于语义层面的组织质量，包括主题的一致性、论述的逻辑性以及段落之间的衔接关系。在医疗大语言模型的安全评估中，连贯性指标可以辅助判断模型输出是否"看起来合理"。许多成功的攻击不仅要求模型生成错误信息，还要求这些信息以看似合理的逻辑组织起来，使其不易被快速识别。如果攻击生成的输出语无伦次、前后矛盾，用户很容易察觉异常；但如果输出逻辑连贯、论证完整，即使内容错误，也可能被误信。因此，连贯性作为攻击隐蔽性的辅助度量，与困惑度形成互补：困惑度衡量表面流畅性，连贯性衡量语义组织性。

连贯性的计算通常基于语义相似度模型。本文采用基于SimCSE（Simple Contrastive Learning of Sentence Embeddings）的度量方法\cite{gao2021simcse}，其数学表达为：
\begin{equation}
\text{Coherence} = \frac{1}{N} \sum_{i=1}^{N} \text{SimCSE}(I_i, O_i)
\label{eq:coherence}
\end{equation}
其中$I_i$是输入指令，$O_i$是模型输出，$N$是测试集大小，$\text{SimCSE}(\cdot)$返回两个句子之间的语义相似度得分。SimCSE是一种基于对比学习的句子嵌入方法，它通过在编码器的dropout噪声之间建立对比关系，学习到高质量的句子表示。使用SimCSE计算连贯性的逻辑是：如果模型的输出与输入指令在语义上保持高度相关，说明输出围绕问题展开，具有较好的连贯性；如果相似度很低，说明输出可能偏离问题主题或存在逻辑跳跃。在医疗问答场景中，高质量的回答应该紧扣患者的问题，针对症状提供诊断或治疗建议，这种相关性应当体现在语义嵌入的相似度上。

连贯性作为安全评估指标需要注意其局限性。首先，高连贯性并不等同于内容正确性。一个完全错误的医疗建议，如果逻辑自洽且围绕问题展开，仍然可以获得较高的连贯性得分。这说明连贯性更适合作为质量指标而非准确性指标。其次，连贯性的测量依赖于参考模型（SimCSE）本身的质量，如果参考模型对医学文本的理解存在偏差，可能导致连贯性评估的不准确。第三，医疗场景中存在一些看似"不连贯"但实际上合理的回答，例如医生指出患者的提问本身存在问题，需要纠正后再回答，这类回答的语义相似度可能偏低，但实际上是恰当的医疗建议。因此，在使用连贯性进行安全评估时，需要结合具体应用场景进行合理解释，避免机械地依赖相似度数值。

在多智能体协作场景下，连贯性指标可以用于评估信息在智能体链中传递时是否保持语义一致性。上游智能体的输出是下游智能体的输入，如果二者之间的语义连贯性过低，说明信息在传递过程中发生了严重的语义偏移或丢失，这可能预示着协作链中的理解错误。通过测量相邻智能体之间输入输出的连贯性，可以定位语义漂移发生的位置。此外，连贯性还可以用于检测智能体是否"忽略"了上游信息而给出无关输出，这在多智能体系统中是一种常见但隐蔽的失效模式。例如，诊断智能体给出了详细的诊断分析，但治疗智能体完全忽略这些分析而给出了不匹配的治疗方案，这种情况下诊断-治疗对之间的连贯性会显著降低。本文的实验将利用连贯性指标分析多智能体协作中的信息传递质量，揭示语义连贯性与错误传播之间的关系。综合ASR、PPL和COH三个指标，可以构建一个多维度的医疗大模型安全评估框架：ASR衡量显性攻击的成功率，PPL衡量攻击的表面隐蔽性，COH衡量输出的语义组织质量，三者结合能够更全面地刻画模型的安全特征。

前述ASR、PPL、COH三个指标均可由自动化系统计算，但自动化评估的可靠性本身需要验证。为此，本文引入人机一致性指标，通过对比自动化评估结果与人类专家判断的吻合程度，验证评估系统的可信度。

\subsection{Cohen's Kappa系数与人机一致性验证}
\label{subsec:kappa}

在医疗大语言模型的安全评估中，一个根本性的挑战是如何验证自动化评估结果的可靠性。当评估系统本身由AI驱动时，存在"AI评估AI"的循环论证风险：如果评估模型与被测模型存在相似的偏差或盲点，自动化评估可能给出错误的高分，掩盖实际的安全隐患。为解决这一问题，本文引入人工专家复核作为金标准，采用Cohen's Kappa系数来衡量自动化评估与人工评估之间的一致性程度。Cohen's Kappa是一种统计指标，用于衡量两个评估者（本场景中为自动化评估Agent和人类专家）对同一组样本进行分类评判时的一致性，其优势在于能够排除随机一致性因素的影响，给出更加客观的一致性评估。

Cohen's Kappa的数学定义为：
\begin{equation}
\kappa = \frac{p_o - p_e}{1 - p_e}
\label{eq:kappa}
\end{equation}
其中$p_o$是观察一致性（Observed Agreement），即评估者之间实际达成一致的样本比例；$p_e$是期望一致性（Expected Agreement），即评估者随机给出评判时预期达成一致的比例\cite{mchugh2012kappa}。Kappa系数的取值范围为$[-1, 1]$，值为1表示完全一致，值为0表示一致性与随机水平相当，值为负表示一致性低于随机水平。在学术界，通常采用Landis和Koch提出的Kappa值解释标准：0.81-1.00表示几乎完全一致（Almost Perfect Agreement），0.61-0.80表示高度一致（Substantial Agreement），0.41-0.60表示中等一致（Moderate Agreement），0.21-0.40表示一般一致（Fair Agreement），0.00-0.20表示轻微一致（Slight Agreement）。对于医疗AI安全评估这一关键应用场景，理想状态下自动化评估与人工评估应达到至少高度一致的水平（Kappa > 0.6），才能认为自动化评估结果具有可信度。

人工专家复核实验的设计需要考虑多个维度。首先，参与复核的专家应当具备相应的医学专业背景，能够准确判断模型输出的医疗建议是否存在错误或不当之处。其次，复核样本的抽样策略应当具有代表性，覆盖不同攻击类型、不同严重程度的错误输出，避免样本偏差影响Kappa计算的可靠性。第三，复核过程需要建立清晰的评判标准和统一的任务描述，确保不同专家的评判具有一致性。在实践中，通常采用多名专家独立评判同一批样本，然后计算专家间的一致性（Inter-Rater Reliability），只有当专家间一致性达到可接受水平时，才将其结果作为金标准来与自动化评估对比。此外，人工复核的成本相对较高，需要在有限的样本预算下实现最大的评估价值，这通常涉及分层抽样等统计学方法，确保各关键类别都有足够的样本被覆盖。

Kappa系数在多智能体安全评估中具有特殊的应用价值。多智能体系统的输出往往经过多个智能体的协同处理，错误的形成和传播机制更加复杂，自动化评估系统的判断可能存在系统性的偏差。通过计算Kappa系数，可以量化自动化评估在哪些场景下与人工判断较为一致，在哪些场景下存在显著差异。如果发现某些类型的安全威胁（如特定的投毒攻击模式）下Kappa值较低，说明自动化评估对该类威胁的检测能力不足，需要针对性地改进评估算法或提示词设计。反之，如果Kappa值整体较高，说明自动化评估系统能够较为可靠地替代人工评估，实现大规模、高效率的安全测试。本文第五章的实验将报告自动化评估Agent与人工专家复核之间的Kappa系数，验证MLSE框架评估结果的可靠性。需要强调的是，即使Kappa值较高，人工专家在以下方面仍不可替代：发现新型攻击模式、识别自动化评估未覆盖的风险边界、提供评估结果的合理解释等。人机协同的评估范式——自动化负责大规模筛查，人工负责深度分析和验证——是医疗AI安全评估的可行路径。

\section{本章小结}
\label{sec:summary}

本章系统阐述了医疗大语言模型与多智能体系统的相关理论基础，为后续的框架设计和实证分析提供了理论支撑。首先介绍了Transformer架构与注意力机制的核心原理，分析了基于该架构的大语言模型在医疗文本处理中的技术优势，进而深入讨论了指令微调技术在医疗领域的应用特点及其安全性影响。从供应链安全视角，本章分析了开源基座模型可能潜伏的后门风险及其通过医疗微调向下游应用传递的机制，指出了模型供应链安全在医疗AI系统中的关键地位。这一分析为理解数据投毒和后门攻击在医疗场景下的特殊性提供了理论基础。

其次，本章对智能体架构进行了形式化定义，提出$\text{Agent} = \langle P, M, A \rangle$的三元组表示方法，明确了感知、记忆、动作三个核心模块在医疗智能体中的功能定位。在此基础上，分析了链式思维推理与医疗诊疗流程的同构性，论证了序列协作模式在多智能体医疗系统中的自然适用性，同时也指出了上游错误在协作链中级联传播的风险。这些理论分析为第三章中MLSE框架的构建和多智能体威胁建模提供了设计依据。

第三，本章系统梳理了四类对抗性攻击机制的原理：数据投毒攻击通过污染训练数据使模型内化错误行为；后门攻击利用触发器机制实现隐蔽的恶意行为激活；提示词注入攻击通过指令覆盖欺骗模型执行非预期操作；越狱攻击针对安全对齐机制试图突破模型的安全约束。通过对各类攻击技术原理的剖析，揭示了医疗大模型面临的多层次安全威胁，为第四章的实证实验设计提供了攻击方法学基础。

第四，本章提出了一个多维度的医疗大模型安全评估指标体系。重新定义了攻击成功率ASR为安全下界代理指标，提出文本困惑度PPL作为攻击隐蔽性指标的新解读，引入连贯性COH衡量输出的语义组织质量，并采用Cohen's Kappa系数验证人机评估一致性。这些指标共同构成了量化评估医疗大模型安全性的工具集，其相互配合能够从不同角度刻画模型的安全特征，为MLSE框架的自动化评估能力提供了度量基础。

综上所述，本章建立了从模型架构、智能体设计、攻击机制到评估指标的完整理论框架。这一框架不仅解释了医疗大模型和多智能体系统如何工作、如何被攻击、如何被评估，更重要的是揭示了各环节之间的内在联系：模型架构决定了其对抗脆弱性，智能体协作模式影响错误传播规律，攻击机制决定了评估指标的设计需求。这些理论认识将在后续章节中转化为具体的框架实现和实证分析，最终形成对医疗多智能体协作系统安全性的全面评估与深入理解。
